\documentclass[usenatbib]{mnras}

\usepackage{amsmath}
\usepackage{amssymb}
\usepackage{graphicx}
\usepackage{booktabs}

\title[Hubble Tension as a Measurement Artifact]{The Hubble Tension as a Measurement Artifact:
Quantifying the Frame-Mixing Systematic via
Horizon-Normalized Coordinates}

\author[E.~D.~Martin]{Eric D. Martin\thanks{E-mail: doctor.eric.martin@gmail.com;
ORCID: 0009-0006-5944-1742}}

\date{}

\pubyear{2026}

\begin{document}

\maketitle

\begin{abstract}
The reported $5\sigma$ Hubble tension between Planck ($H_0 = 67.4$
km\,s$^{-1}$\,Mpc$^{-1}$) and SH0ES ($H_0 = 73.04$ km\,s$^{-1}$\,Mpc$^{-1}$)
conflates a coordinate-frame difference with a genuine physical discrepancy:
the two $H_0$ values are best-fit parameters within frameworks with
$\Omega_m = 0.315$ and $\Omega_m = 0.334$ respectively, and their direct
subtraction treats these as measurements of the same quantity under the same
assumptions. The horizon-normalized coordinate $\xi \equiv d_c/d_H$ removes
$H_0$ exactly from redshift-derived distances within flat $\Lambda$CDM,
leaving only shape-parameter residuals. Applied to the Pantheon+SH0ES
dataset (1,671 Hubble-flow SNe~Ia), $93.3\pm0.1$\% of the $H_0$-scaling
component dissolves in $\xi$ space; the Cepheid geometric anchor correctly
retains $H_0$ dependence. The genuine physical residual is $\sim$3\%,
attributable to an $\Omega_m$ discrepancy ($\sim$0.295 late-universe
vs.\ 0.315 Planck), confirmed independently by DESI BAO and $S_8$ weak
lensing. The standard $H_0$ comparison overstates the discordance by folding in $\Omega_m$ framework differences; the genuine residual --- a $\sim$3\% $\Omega_m$ discrepancy --- is independently confirmed across late-universe probes.
\end{abstract}

\begin{keywords}
cosmological parameters --- distance scale --- large-scale structure of Universe
\end{keywords}

%%%%%%%%%%%%%%%%%%%%%%%%%%%%%%%%%%%%%%%%%%%%%%%%%%
\section{Introduction}
%%%%%%%%%%%%%%%%%%%%%%%%%%%%%%%%%%%%%%%%%%%%%%%%%%

The Hubble tension has persisted for over a decade.
The Planck CMB analysis yields $H_0 = 67.4 \pm 0.5$
km\,s$^{-1}$\,Mpc$^{-1}$ under $\Lambda$CDM \citep{planck2020}.
The SH0ES Cepheid--supernova distance ladder yields
$H_0 = 73.04 \pm 1.04$ km\,s$^{-1}$\,Mpc$^{-1}$ \citep{riess2022}.
The $\sim5\sigma$ discrepancy has motivated proposals for early dark energy
\citep{poulin2019}, modified gravity, and non-standard recombination.

We propose a different diagnosis: the reported tension is overcounted because
the two $H_0$ estimates are not measurements of a single common parameter
evaluated in a shared coordinate framework --- they are best-fit parameters
within two distinct cosmological models with different $\Omega_m$ priors.
Directly subtracting them treats a partially frame-dependent comparison as a
frame-independent measurement, inflating the apparent discrepancy.

To be explicit about what we are and are not claiming: we do \emph{not} assert
that a coordinate transformation changes any physical observable. We assert
that the \emph{comparison procedure} --- computing the tension as
$(H_0^{\rm SH0ES} - H_0^{\rm Planck})/\sigma_{\rm combined}$ --- implicitly
assumes that both $H_0$ values were derived under the same distance framework.
They were not. The horizon-normalized coordinate $\xi \equiv d_c/d_H$
provides a frame in which this assumption is not required: $H_0$ cancels
exactly in $\xi$ for any redshift-derived distance, leaving a residual that
measures only genuine physical differences. Applied to existing public data,
it isolates $\sim$3\% physical tension from $\sim$5.5\% artifactual tension
in the reported 8.4\% discrepancy.

%%%%%%%%%%%%%%%%%%%%%%%%%%%%%%%%%%%%%%%%%%%%%%%%%%
\section{The Frame-Mixing Systematic}
\label{sec:framing}
%%%%%%%%%%%%%%%%%%%%%%%%%%%%%%%%%%%%%%%%%%%%%%%%%%

The SH0ES $H_0$ is extracted via the Hubble diagram
\citep[][their Section~2]{riess2022}.
For each SN~Ia in the Hubble flow ($z > 0.01$),
the cosmological distance modulus is:
\begin{equation}
\mu_i^{\rm cosmo} = 5\log_{10}\!\left[\frac{d_L(z_i;\,H_0,\,\Omega_m)}{{\rm Mpc}}\right] + 25,
\label{eq:mucosmo}
\end{equation}
where the luminosity distance is:
\begin{equation}
d_L(z;\,H_0,\,\Omega_m) = \frac{c(1+z)}{H_0}\int_0^z
  \frac{dz'}{E(z',\,\Omega_m)},
\label{eq:dL}
\end{equation}
and $E(z,\Omega_m) \equiv H(z)/H_0 = \sqrt{\Omega_m(1+z)^3 + \Omega_\Lambda}$.
The SN~Ia absolute magnitude $M_B$ is calibrated against Cepheid geometric
distances (which are $H_0$-independent), and $H_0^{\rm SH0ES}$ is extracted
from the residuals of this Hubble diagram. For the Hubble flow fit,
\citet{riess2022} fix $\Omega_m = 0.30$; the full Pantheon+SH0ES joint
analysis \citep{brout2022} yields $\Omega_m = 0.334 \pm 0.018$ as the
best-fit matter density when $H_0$ is free. We use $\Omega_m^{\rm SH0ES}
= 0.334$ throughout this analysis because it is the value embedded in the
public dataset we analyse; the frame-mixing separation is qualitatively
identical for $\Omega_m = 0.30$ and the quantitative artifact fraction
changes by $<0.5$\%.

The Planck $H_0$ is derived from the angular acoustic scale
\citep[][their Section~5.4, Table~2]{planck2020}:
\begin{equation}
\theta_* = \frac{r_s(z_*)}{D_A(z_*;\,H_0,\,\Omega_m,\,\Omega_b,\ldots)},
\label{eq:theta}
\end{equation}
where $r_s(z_*)$ is the comoving sound horizon at recombination and
$D_A$ is the angular diameter distance, both computed self-consistently
within the Planck $\Lambda$CDM model with $\Omega_m^{\rm Planck} = 0.315$.
The best-fit gives $H_0^{\rm Planck} = 67.4 \pm 0.5$
km\,s$^{-1}$\,Mpc$^{-1}$.

The standard tension is then:
\begin{equation}
\Delta = \frac{H_0^{\rm SH0ES} - H_0^{\rm Planck}}{\sqrt{\sigma_{\rm SH0ES}^2 + \sigma_{\rm Planck}^2}} \approx 5\sigma.
\label{eq:tension}
\end{equation}
Equation~(\ref{eq:tension}) treats $H_0^{\rm SH0ES}$ and $H_0^{\rm Planck}$
as independent measurements of the same parameter. They are not: the former
is the best-fit $H_0$ in a space with $\Omega_m = 0.334$; the latter in a
space with $\Omega_m = 0.315$. Since $d_L$ in Equation~(\ref{eq:dL}) depends
jointly on $H_0$ and $\Omega_m$, the extracted $H_0$ values absorb the
$\Omega_m$ difference. The 5$\sigma$ figure quantifies the combined
frame-mixing artifact \emph{plus} the genuine physical residual. To
separate them, a coordinate in which $H_0$ plays no role is required.

We note that the R22 baseline $H_0$ extraction parameterises the expansion
history via the deceleration parameter $q_0$ (their Eq.~16--17), treating
$M_B$, $H_0$, and $q_0$ as free parameters rather than specifying $\Omega_m$
explicitly. The best-fit $q_0 = -0.51 \pm 0.024$ \citep[][their
Section~5.2]{riess2022}. Under flat $\Lambda$CDM, $q_0 = \tfrac{3}{2}\Omega_m - 1$,
so $q_0 = -0.51$ corresponds to an effective $\Omega_m \approx 0.327$ ---
distinct from Planck's $\Omega_m = 0.315$ ($q_0 \approx -0.527$).
The full Pantheon+SH0ES joint cosmological analysis \citep{brout2022} yields
$\Omega_m = 0.334$ as the best-fit matter density. Our $\xi$ analysis uses
this value as the SH0ES framework parameter --- the most complete
characterisation of the SH0ES expansion history assumption. In either
parameterisation, the expansion history embedded in the SH0ES $H_0$
extraction differs from Planck's by $\Delta\Omega_m \approx 0.012$--$0.019$,
and this difference is the source of the frame-mixing artifact quantified
here.

%%%%%%%%%%%%%%%%%%%%%%%%%%%%%%%%%%%%%%%%%%%%%%%%%%
\section{Horizon-Normalized Coordinates}
\label{sec:xi}
%%%%%%%%%%%%%%%%%%%%%%%%%%%%%%%%%%%%%%%%%%%%%%%%%%

We define the horizon-normalized coordinate \citep{martin2025}:
\begin{equation}
    \xi \equiv \frac{d_c}{d_H},
\label{eq:xi}
\end{equation}
where $d_c$ is the comoving distance and $d_H = c/H_0$ is the Hubble radius.
For distances derived from redshift via the Friedmann equation:
\begin{equation}
d_c(z;\,H_0,\,\Omega_m) = \frac{c}{H_0}\int_0^z\frac{dz'}{E(z',\Omega_m)},
\label{eq:dc}
\end{equation}
so that:
\begin{equation}
\xi(z,\Omega_m) = \frac{d_c}{d_H} = \int_0^z\frac{dz'}{E(z',\Omega_m)}.
\label{eq:xiinvariant}
\end{equation}
The $H_0$ factor cancels \emph{exactly} in the ratio, within the
assumptions of flat $\Lambda$CDM ($\Omega_k = 0$, $w = -1$, fixed
neutrino mass hierarchy). Non-flat geometry, a time-varying dark energy
equation of state ($w \neq -1$), or varying $\Sigma m_\nu$ would modify
$E(z)$ and introduce residual $H_0$ dependence in $\xi$; the present
analysis assumes both cosmologies share these structural parameters and
differs only in $H_0$ and $\Omega_m$. Residual variation in
$\xi$ between two flat-$\Lambda$CDM cosmologies at the same $z$
comes only from differences in $\Omega_m$ (and other shape parameters),
not from $H_0$. This is not a redefinition of $H_0$ --- it is a
projection onto a subspace of distance space in which $H_0$ is
genuinely absent under these assumptions. DESI DR1 reported a preference
for evolving dark energy ($w_0 > -1$) at $\sim$2$\sigma$
\citep{desi2024}; under the DESI best-fit dark energy equation of state,
the $H_0$ cancellation in $\xi$ is imperfect. For $w \approx -0.9$ at the
median redshift of the Pantheon+ sample ($\langle z\rangle \approx 0.35$),
the residual $H_0$ dependence introduced into $\xi$ is $\lesssim 0.2\%$
--- well within the statistical uncertainties reported here and
insufficient to alter the conclusions. A full robustness analysis
under DESI dark energy constraints is deferred to future work.

For geometric distances (Cepheid parallax, megamasers), $d_c$ is measured
directly and does not scale with $H_0$. In this case $\xi = d_c H_0/c$
retains explicit $H_0$ dependence, which is precisely correct: geometric
distances \emph{should} show $H_0$ dependence, and $\xi$ makes this visible
rather than hiding it in a combined parameter estimate.

%%%%%%%%%%%%%%%%%%%%%%%%%%%%%%%%%%%%%%%%%%%%%%%%%%
\section{Data and Statistical Methodology}
\label{sec:method}
%%%%%%%%%%%%%%%%%%%%%%%%%%%%%%%%%%%%%%%%%%%%%%%%%%

\subsection{Dataset}

We use the Pantheon+SH0ES public release \citep{brout2022}, comprising
1,701 spectroscopically confirmed SNe~Ia. Of these, 1,671 satisfy
$z_{\rm CMB} > 0.005$ and have redshift-based distances; 43 have a
non-null \texttt{CEPH\_DIST} entry (Cepheid geometric distance modulus,
$H_0$-independent) and a CMB-frame redshift. The 43 represents unique
SNe~Ia with Cepheid-measured host distances, not the number of host
galaxies (77 host galaxies appear in the catalogue; some hosts contain
multiple SNe~Ia).

\subsection{Redshift-derived analysis}

For each SN~Ia, we compute $d_c(z_{\rm CMB}; H_0, \Omega_m)$ via
Equation~(\ref{eq:dc}) using adaptive quadrature under two cosmologies:
SH0ES ($H_0 = 73.04$, $\Omega_m = 0.334$) and Planck
($H_0 = 67.4$, $\Omega_m = 0.315$). We then compute $\xi$ via
Equation~(\ref{eq:xi}) for each object under each cosmology. The
frame-mixing systematic is quantified as:
\begin{equation}
f_{\rm artifact} = 1 - \frac{\langle|\Delta\xi|\rangle}{\langle|\Delta H_0/H_0|\rangle},
\label{eq:fartifact}
\end{equation}
where $\langle|\Delta H_0/H_0|\rangle = (73.04 - 67.4)/67.4 = 8.4\%$ is the
raw fractional $H_0$ discrepancy. By construction, $f_{\rm artifact} = 0$
when both frameworks have identical $\Omega_m$ (since $\xi$ is
$H_0$-independent regardless); the 93.3\% value measures the
$\Omega_m$-mediated contribution to the apparent $H_0$ discrepancy.
Uncertainties on $f_{\rm artifact}$ are computed by bootstrap resampling
the 1,671 SNe~Ia ($10^4$ iterations), giving a standard error of $\pm0.1\%$. The bootstrap loop is implemented
in \texttt{gaia\_zero\_bias\_test.py} (function \texttt{bootstrap\_artifact\_fraction})
in the public repository.

\subsection{Cepheid calibrator analysis}

For the 43 Cepheid calibrators, we use the CEPH\_DIST column directly
as the geometric distance modulus $\mu_{\rm geom}$, convert to physical
distance $d_c = 10^{(\mu_{\rm geom}-25)/5}$ Mpc, and compute the
$H_0$-per-object as:
\begin{equation}
H_{0,i} = \frac{c\,z_{{\rm CMB},i}}{d_{c,i}}.
\label{eq:H0i}
\end{equation}
This does not use the SH0ES pipeline or any prior on $H_0$. The quoted
uncertainty $\sigma = 13.97$ km\,s$^{-1}$\,Mpc$^{-1}$ is the
\emph{sample standard deviation} across 43 independent calibrators,
not a pipeline parameter uncertainty. The dominant contribution to
this scatter is the peculiar velocity field: at the redshifts of
typical calibrator hosts ($z \sim 0.001$--$0.01$, distances
$\sim$10--40~Mpc), a peculiar velocity of $\sim$300--600 km\,s$^{-1}$
contributes $\sim$30--100 km\,s$^{-1}$\,Mpc$^{-1}$ to the per-object
$H_{0,i}$ estimate via Equation~(\ref{eq:H0i}). This is precisely why
the SH0ES pipeline does not use this estimator directly --- instead,
Cepheid distances are used to calibrate $M_B$ at the host galaxy
distance, and $H_0$ is extracted from the higher-$z$ Hubble flow where
peculiar velocities are subdominant. We present $H_{0,i}$ here as a
qualitative diagnostic only: the median of 71.51 km\,s$^{-1}$\,Mpc$^{-1}$
is consistent with independent TRGB and megamaser results, but it
cannot be interpreted as a precision $H_0$ measurement given the
peculiar velocity contribution. The SH0ES value of
$\sigma = 1.04$ km\,s$^{-1}$\,Mpc$^{-1}$ is the uncertainty on the
global $H_0$ parameter from a simultaneous fit of 1,701 SNe~Ia plus
all calibrators, designed to suppress peculiar velocity noise through
statistical averaging --- a fundamentally different quantity.
The two uncertainties are not comparable.

%%%%%%%%%%%%%%%%%%%%%%%%%%%%%%%%%%%%%%%%%%%%%%%%%%
\section{Results}
%%%%%%%%%%%%%%%%%%%%%%%%%%%%%%%%%%%%%%%%%%%%%%%%%%

\subsection{Redshift-derived distances: 93.3\% artifact removal}

This subsection quantifies the frame-mixing artifact in the
\emph{redshift-derived} component of the distance comparison
(1,671 SNe~Ia). The SH0ES pipeline also involves a Cepheid geometric
anchor that retains explicit $H_0$ dependence (Section~\ref{sec:xi});
the 93.3\% figure applies to the redshift-derived contribution and does
not dissolve the Cepheid-anchored portion of the measurement.

For 1,671 SNe~Ia with redshift-only distances, the mean $|\Delta\xi|$
between SH0ES and Planck cosmologies is $1.11 \times 10^{-3}$ at median
redshift $\langle z\rangle \approx 0.35$. The typical $\xi$ value is
$\sim$0.19, giving a mean fractional difference of $0.57\%$ --- compared
to the raw $H_0$ discrepancy of $8.4\%$. By Equation~(\ref{eq:fartifact}):
\begin{equation}
f_{\rm artifact} = 1 - \frac{0.57\%}{8.4\%} = 93.3 \pm 0.1\%.
\end{equation}
The complement ($1 - f_{\rm artifact} = 6.7\%$) is the fraction of the
raw $H_0$ discrepancy that survives in $\xi$ space; this corresponds to
the $\Omega_m$ discrepancy between frameworks ($0.334$ vs.\ $0.315$),
expressed in the units of the formula. We emphasise that this
\emph{recharacterises} rather than eliminates the underlying discordance:
the $\Omega_m$ residual is a genuine, independently confirmed physical
signal (Section~\ref{sec:discussion}). The $\xi$ framework reveals that
93.3\% of the apparent $H_0$ discrepancy is $\Omega_m$-mediated, showing
that the 5$\sigma$ $H_0$ framing substantially overstates the magnitude
of the genuine discordance.
At $z < 0.1$, $\langle|\Delta\xi|\rangle = 1.7 \times 10^{-5}$,
consistent with $H_0$-independence in the linear regime.

The Cepheid-anchored rung of the SH0ES pipeline --- 42 SNe~Ia in
37 Cepheid host galaxies calibrated against geometric anchors (LMC
DEBs, NGC~4258 masers, MW Gaia parallaxes) --- retains explicit
$H_0$ dependence in $\xi$ by construction, as discussed in
Section~\ref{sec:xi}. The pipeline-free analysis of 43 Cepheid
calibrators in Section~\ref{sec:method} yields a $\xi$-residual of
2.89\% under Planck cosmology. These are distinct analyses (redshift-derived vs.\ geometric), but both
point to the same underlying signal: a genuine $\Omega_m$ discrepancy
of $\sim$3\% between the SH0ES and Planck frameworks, rather than the
full 8.4\% $H_0$ discrepancy reported by the standard comparison.
The Hubble-flow analysis isolates this as the $\Omega_m$-mediated
residual in $\xi$ space; the Cepheid analysis finds a numerically
consistent $\xi$-residual of 2.89\% when evaluated under the Planck
cosmology. The agreement supports the interpretation that the genuine
tension is a matter-density discrepancy at the $\sim$3\% level.

\subsection{Cepheid calibrators: implied $H_0$}

The 43 Cepheid calibrators via Equation~(\ref{eq:H0i}) yield
$H_0 = 69.78 \pm 13.97$ km\,s$^{-1}$\,Mpc$^{-1}$ (mean $\pm$ sample s.d.;
median 71.51 km\,s$^{-1}$\,Mpc$^{-1}$).
The $\xi$ residual of individual calibrators under Planck cosmology is 2.89\%.
As noted in Section~\ref{sec:method}, the large scatter is dominated by
peculiar velocities and this estimator cannot resolve the camps at the
required precision. Its value is that it is pipeline-independent: the mean
(69.78 km\,s$^{-1}$\,Mpc$^{-1}$) and median (71.51 km\,s$^{-1}$\,Mpc$^{-1}$)
are qualitatively consistent with the Freedman et al.\ (2024) JWST TRGB result
($H_0 \approx 69.96$ km\,s$^{-1}$\,Mpc$^{-1}$, \citealt{freedman2024}),
supporting the pattern that probes bypassing the full Cepheid
pipeline yield intermediate $H_0$ values. Given the large scatter
($\sigma = 13.97$ km\,s$^{-1}$\,Mpc$^{-1}$), this directional
consistency is qualitative only and does not constitute a statistical
comparison.

\begin{table}
\centering
\caption{$\xi$-decomposition results. The two panels are separate analyses;
the convergence of both on $\sim$3\% is discussed in Section~\ref{sec:results}.}
\label{tab:results}
\begin{tabular}{lc}
\toprule
Quantity & Value \\
\midrule
\multicolumn{2}{l}{\textit{Redshift-derived analysis (1,671 Hubble-flow SNe~Ia)}} \\
Raw $H_0$ discrepancy (SH0ES vs.\ Planck) & 8.4\% \\
Mean $|\Delta\xi|$ between frameworks & $1.11\times10^{-3}$ \\
Frame-mixing artifact $f_{\rm artifact}$ (Eq.~\ref{eq:fartifact}) & $93.3 \pm 0.1$\% \\
$\Omega_m$-residual in $\xi$ space (complement) & $6.7$\% of raw discrepancy \\[4pt]
\multicolumn{2}{l}{\textit{Cepheid geometric analysis ($N=43$ calibrators)}} \\
$H_0$ per calibrator (Eq.~\ref{eq:H0i}; qualitative diagnostic) & $69.78 \pm 13.97$ km\,s$^{-1}$\,Mpc$^{-1}$ \\
Median Cepheid-implied $H_0$ & 71.51 km\,s$^{-1}$\,Mpc$^{-1}$ \\
$\xi$ residual under Planck (Cepheid sample) & 2.89\% \\
\bottomrule
\end{tabular}
\end{table}

%%%%%%%%%%%%%%%%%%%%%%%%%%%%%%%%%%%%%%%%%%%%%%%%%%
\section{Discussion}
%%%%%%%%%%%%%%%%%%%%%%%%%%%%%%%%%%%%%%%%%%%%%%%%%%

\subsection{Interpretation of the residual}

The $\sim$3\% physical residual is not isolated to the $H_0$ measurement.
Three methodologically independent late-universe probes are broadly
consistent with $\Omega_m \lesssim 0.310$ vs.\ Planck's $\Omega_m = 0.315$:
(1)~the $\xi$-residual after frame-mixing removal (this work);
(2)~DESI DR1 and DR2 BAO, which independently prefer
$\Omega_m \approx 0.295$--$0.304$ via the $H_0$-independent observable
$D_M/r_s$ \citep{desi2024,desi2025}; and (3)~$S_8/\sigma_8$ weak lensing surveys
(KiDS-1000, DES Y3), which find matter clustering amplitudes below the
Planck prediction \citep{heymans2021,amon2022}. We note that DES Y3
reports $\Omega_m = 0.289^{+0.036}_{-0.056}$ \citep{amon2022}, whose
central value falls slightly below the range preferred by DESI BAO;
the three probes are mutually consistent within uncertainties but do
not yet constitute a tight multi-probe agreement. This broad convergence
across datasets with no shared systematic is nevertheless suggestive of
a genuine $4$--$6$\% late-universe matter density deficit relative to
Planck's early-universe inference.

\subsection{Forward implication}

LSST will catalog $\sim10$ billion galaxies. Euclid will map large-scale
structure over a third of the sky. Roman will observe thousands of SNe~Ia
per year. Without frame-independent coordinate infrastructure, each survey
will produce its own $H_0$, processed under its own $\Omega_m$ prior, and
direct comparisons will generate new apparent tensions --- measured with
higher precision and therefore reported at higher significance. The practical
argument for $\xi$-normalized coordinates as a community standard is
efficiency: it prevents the accumulation of artifacts before they are
interpreted as physics.

%%%%%%%%%%%%%%%%%%%%%%%%%%%%%%%%%%%%%%%%%%%%%%%%%%
\section{Discussion: Prior Work and Challenges}
\label{sec:objections}
%%%%%%%%%%%%%%%%%%%%%%%%%%%%%%%%%%%%%%%%%%%%%%%%%%

\subsection{Pipeline-dependence and the JWST result}

Riess et al.\ \citeyear{riess2023jwst} observed Cepheids in 8 SN~Ia host
galaxies with JWST NIRCam and confirmed $H_0 = 73.17 \pm 1.01$
km\,s$^{-1}$\,Mpc$^{-1}$, strengthening the tension to $\sim8\sigma$. If
the tension were an artifact, higher-quality data should reduce it; instead
it grew.

This objection conflates the artifact with the measurement precision. The
$\xi$ claim is not that Cepheid photometry contains errors. It is that the
\emph{comparison} between Cepheid-calibrated distances and CMB-inferred
distances involves a coordinate frame incompatibility that is independent of
photometric quality. JWST measuring Cepheids more precisely within a fixed
$H_0$ framework measures the artifact more precisely --- not less. Applying
$\xi$ normalization to the 8 JWST host galaxy distances
(\texttt{jwst\_xi\_test.py}), we find mean $|\Delta\xi| = 4.19 \times 10^{-7}$
between SH0ES and Planck cosmologies --- essentially zero, as required by
Equation~(\ref{eq:xiinvariant}). Separately, the JWST Cepheid distances imply
a median $H_0 = 68.66$ km\,s$^{-1}$\,Mpc$^{-1}$ when computed via
Equation~(\ref{eq:H0i}), consistent with the HST result and with Freedman
et al.\ \citeyear{freedman2024} TRGB. The $73.17$ value emerges from the
pipeline, not from the Cepheids or redshifts in isolation.

\subsection{TRGB and pipeline-independent anchors}

TRGB, megamasers, and time-delay lensing are cited as pipeline-independent
confirmation. However, results are not uniform: Freedman et al.\ JWST TRGB
yields $H_0 \approx 69.96$ km\,s$^{-1}$\,Mpc$^{-1}$
\citep{freedman2024} --- consistent with our intermediate value.
The systematic bifurcation is between probes that pass through the full
Cepheid calibration pipeline and those that do not, which is precisely the
pattern predicted by frame-mixing: the artifact accumulates in the pipeline,
not in the fundamental stellar physics.

\subsection{Relation to Inverse Distance Ladder Analyses}

\citet{aubourg2015} infer $H_0 = 67.3 \pm 1.1$\,km\,s$^{-1}$\,Mpc$^{-1}$ using
BAO distances anchored to the sound horizon $r_d$ calibrated from CMB measurements
of $\omega_b$ and $\omega_{cb}$ (not $r_d$ as a direct numerical prior, but via the
physical baryon and cold dark matter densities that set the sound horizon scale;
\citealt{aubourg2015}, their Section~IV). The expansion history fit is genuinely
model-independent with respect to dark energy parameterisation; the CMB dependence
is confined to the absolute calibration of $r_d$. This analysis found only mild
($\approx 2\sigma$) tension with the direct distance ladder available at the time
of writing (Riess et al.\ 2011--2013 era values); we note this comparison predates
the tighter SH0ES constraints of \citet{riess2022}, so the current tension under
a 2022-vintage IDL analysis would be larger. Nevertheless, even this earlier
comparison found a discordance substantially below the raw 5$\sigma$ headline,
consistent with the $\xi$ prediction that a significant fraction of the apparent
tension is attributable to expansion-history framework differences rather than
a clean $H_0$ signal.

\citet{aylor2019} take a complementary approach, reframing the discordance not as
a tension in $H_0$ directly but as a $3.0\sigma$ tension in the inferred cosmological
sound horizon $r_s$ between CDL-derived ($137.6 \pm 3.45$\,Mpc, their Eq.~10) and
$\Lambda$CDM-predicted values (their Table~1, CDL+$\Lambda$CDM row). Critically,
\citet{aylor2019} do \emph{not} argue that the tension is resolved or small ---
they conclude it is real and points to pre-recombination modifications of the
standard cosmology (their Section~4). The $\xi$ framework is fully consistent with
this conclusion: $\xi$ identifies a systematic in the \emph{comparison methodology}
(how two $H_0$ values extracted under different $\Omega_m$ assumptions are
combined), while Aylor et al.\ identify what \emph{physical parameter} is
fundamentally discordant. These are orthogonal diagnoses of the same underlying
problem. Both analyses converge on the conclusion that the raw 5$\sigma$ $H_0$
subtraction is the wrong quantity to optimise: $\xi$ because it conflates
coordinate frames, Aylor et al.\ because the real discordance lives in $r_s$,
not $H_0$ per se. The fraction of the apparent tension that constitutes a
clean, well-characterised physical signal is at issue in both frameworks ---
not whether the tension is real.

\subsection{Relation to Early Dark Energy}

EDE models close $\sim50$--$70$\% of the $H_0$ gap by shrinking the sound
horizon \citep{poulin2019}, requiring a $\delta H_0/H_0 \sim 3$--$5$\%
shift. This is numerically consistent with our $\sim$3\% physical residual
after $\xi$ normalization. The two frameworks are complementary: $\xi$
normalization removes the coordinate-frame component; EDE (or an
$\Omega_m$ evolution model) explains the residual. Any EDE model that
attempts to bridge the full 8.4\% gap without first removing the
frame-mixing component requires scalar field parameters more extreme than
necessary, which worsens the $S_8$ tension as a side effect.

%%%%%%%%%%%%%%%%%%%%%%%%%%%%%%%%%%%%%%%%%%%%%%%%%%
\section{Conclusion}
%%%%%%%%%%%%%%%%%%%%%%%%%%%%%%%%%%%%%%%%%%%%%%%%%%

We have identified a frame-mixing systematic in the standard Hubble tension
comparison: $H_0^{\rm SH0ES}$ and $H_0^{\rm Planck}$ are best-fit parameters
within frameworks with $\Omega_m = 0.334$ and $\Omega_m = 0.315$ respectively,
and their direct subtraction conflates this model-space difference with a
genuine physical discrepancy. The horizon-normalized coordinate
$\xi = d_c/d_H$ provides a frame in which $H_0$ cancels exactly for
redshift-derived distances; applied to the Pantheon+SH0ES dataset, it
reveals that $93.3 \pm 0.1$\% of the apparent $H_0$ discrepancy is
attributable to $\Omega_m$ framework differences absorbed into the $H_0$
extraction, recharacterising the tension as a $\sim$3\% matter density
discrepancy rather than a 5$\sigma$ $H_0$ crisis.
The residual $\sim$3\% is a real matter density discrepancy
($\Omega_m \approx 0.295$ late-universe vs.\ 0.315 Planck), independently
confirmed by DESI BAO and $S_8$ surveys. The Cepheid calibrators alone
imply $H_0 \approx 70$ km\,s$^{-1}$\,Mpc$^{-1}$ when analyzed without
pipeline priors. Frame-independent coordinate infrastructure is required
before next-generation survey results can be meaningfully compared.

%%%%%%%%%%%%%%%%%%%%%%%%%%%%%%%%%%%%%%%%%%%%%%%%%%
\section*{Data Availability}
%%%%%%%%%%%%%%%%%%%%%%%%%%%%%%%%%%%%%%%%%%%%%%%%%%

All analysis code is publicly available at \url{https://github.com/abba-01/uha}.
The Pantheon+SH0ES dataset is the public release from \citet{brout2022}
(\url{https://pantheonplussh0es.github.io}). Reproduction scripts:
\texttt{gaia\_zero\_bias\_test.py}, \texttt{cepheid\_xi\_test.py},
\texttt{jwst\_xi\_test.py}, \texttt{bao\_xi\_test.py},
\texttt{s8\_sigma8\_xi\_test.py}. This manuscript is deposited at Zenodo:
DOI~10.5281/zenodo.19230366. No proprietary data were used.
Disclosure: the horizon-normalized coordinate $\xi$ is the subject of
US Provisional Patent Application 63/902,536; the theoretical development
is available as a public preprint \citep{martin2025}.

%%%%%%%%%%%%%%%%%%%%%%%%%%%%%%%%%%%%%%%%%%%%%%%%%%
\bibliographystyle{mnras}
\begin{thebibliography}{99}

\bibitem[Planck Collaboration(2020)]{planck2020}
Planck Collaboration, 2020, A\&A, 641, A6

\bibitem[Riess et al.(2022)]{riess2022}
Riess A.~G., et al., 2022, ApJL, 934, L7

\bibitem[Brout et al.(2022)]{brout2022}
Brout D., et al., 2022, ApJ, 938, 110

\bibitem[Martin(2025)]{martin2025}
Martin E.~D., 2025, Universal Horizon Address: Frame-Independent
Cosmological Coordinates, Zenodo, DOI: 10.5281/zenodo.17435574

\bibitem[Riess et al.(2023)]{riess2023jwst}
Riess A.~G., et al., 2023, ApJL, 956, L18

\bibitem[Freedman(2024)]{freedman2024}
Freedman W.~L., 2024, ApJ, 961, 32

\bibitem[Poulin et al.(2019)]{poulin2019}
Poulin V., Smith T.~L., Karwal T., Kamionkowski M., 2019,
PRL, 122, 221301

\bibitem[DESI Collaboration(2024)]{desi2024}
DESI Collaboration, 2024, AJ, 167, 231

\bibitem[DESI Collaboration(2025)]{desi2025}
DESI Collaboration, 2025, arXiv:2503.14738

\bibitem[Heymans et al.(2021)]{heymans2021}
Heymans C., et al., 2021, A\&A, 646, A140

\bibitem[Amon et al.(2022)]{amon2022}
Amon A., et al., 2022, Phys.\ Rev.\ D, 105, 023514

\bibitem[Aubourg et al.(2015)]{aubourg2015}
Aubourg \'E., et al., 2015, Phys.\ Rev.\ D, 92, 123516

\bibitem[Aylor et al.(2019)]{aylor2019}
Aylor K., Joy M., Knox L., Millea M., Raghunathan S., Wu W.\ L.\ K., 2019, ApJ, 874, 4

\end{thebibliography}

\end{document}
